\documentclass{situationist-preprint}

\title{Meaning, Externalism and Internal States}
\msid{LNG-08-a}
\msfamily{LNG-08}
\msdiscipline{Philosophy of Language}
\mstype{Principal manuscript}
\mskeywords{semantic externalism; Twin Earth; narrow content; natural kind terms; two-dimensionalism; self-knowledge}

\begin{document}
\maketitle

\begin{abstract}
Are meanings in the head? Putnam's Twin Earth argument shows that two speakers in the same internal state can mean different things by \enquote{water}, because the extension of the word depends on the nature of the samples it was grounded in. We accept the argument: meanings, understood as what determines extension, are not in the head. But they are not independent of the head either. We argue that the meaning of a term is fixed jointly by the environment and by three things that are in the head: the sortal under which its grounding samples are taken, the speaker's deferential and sameness intentions, and a body of information that serves as a reference-fixing reserve. None of these is a description that the referent must satisfy when the grounding succeeds, which is why kind terms refer independently of the descriptions associated with them; the information becomes relevant only when the causal grounding fails, as the theory of reference of LNG-02 predicts. We show that this division answers the narrow-content theorist, the two-dimensionalist and the argument from self-knowledge.
\end{abstract}

\section{The problem}

\citet{Putnam1975} argued that two assumptions of traditional semantics cannot both be true: that knowing the meaning of a term is a matter of being in a certain psychological state, and that the meaning of a term determines its extension. His argument was the Twin Earth case. Twin Earth is like Earth except that the liquid that fills its lakes and rivers, falls as rain and is drunk by its inhabitants is not H$_2$O but a different substance, XYZ, indistinguishable from water in ordinary conditions. In 1750, before chemistry could distinguish them, Oscar on Earth and his molecule-for-molecule duplicate on Twin Earth were in the same internal state; yet Oscar's word \enquote{water} applied to H$_2$O and not to XYZ, and his twin's applied to XYZ and not to H$_2$O. Since the meaning of the word determines its extension, the meanings differ; since the internal states are the same, meanings are not determined by internal states. In Putnam's slogan, meanings just ain't in the head.

\citet{Burge1979} extended the argument from natural kind terms to terms in general. A patient who believes that he has arthritis in his thigh, not knowing that arthritis is by definition an inflammation of the joints, has a false belief about arthritis; his internal duplicate, in a community where \enquote{arthritis} is used for a wider class of rheumatoid ailments, has a true belief about something else. What a speaker's words mean depends on the practices of the community to which she defers.

Two responses have been influential. Internalists have proposed that, besides the externally individuated meanings Putnam identified, there is a narrow content shared by Oscar and his twin, which is in the head and which is what psychology needs \citep{Fodor1987,Segal2000}. Two-dimensionalists have proposed that a term has two intensions, one of which, the primary or epistemic intension, is fixed by the speaker's internal state and captures how the term picks out its referent in whatever world turns out to be actual \citep{Jackson1998,Chalmers2002}. Both responses hold that something meaning-like is in the head. Neither denies Putnam's conclusion about the extension-determining meaning.

We defend a thesis that sits between Putnam's slogan and the internalist responses. Meanings are not in the head. But they depend on what is in the head, in a precise sense: the meaning of a term is fixed by the speaker's internal state together with her environment and community, and it varies with the internal state when the environment is held fixed. The contribution of the head is not a description that the referent must satisfy. It consists of the conditions that make a grounding a grounding of one thing rather than another, and of a reserve that fixes reference when the grounding fails. Section~2 states the framework, Section~3 proves the two halves of the thesis, Section~4 explains how the head's contribution can be indispensable and still non-descriptive, and Section~5 answers objections.

\section{Framework}

\begin{definition}\label{d:meaning}
The \emph{meaning} of a term, as used by a speaker, is whatever determines its extension, in the actual world and in counterfactual situations.
\end{definition}

Definition~\ref{d:meaning} is Putnam's second assumption, and we adopt it. There may be other things worth calling meaning, such as the cognitive significance that distinguishes \enquote{Hesperus} from \enquote{Phosphorus}; Section~5 returns to them.

\begin{definition}\label{d:internal}
A speaker's \emph{internal state} is the totality of her intrinsic physical states and of the psychological states that supervene on them. Two speakers are \emph{internal duplicates} if they are in the same internal state.
\end{definition}

\begin{definition}\label{d:inhead}
Meanings are \emph{in the head} if internal duplicates mean the same by every term they share. Meanings \emph{depend on the head} if there are speakers who share their environment and their community but differ in internal state, and who differ in what they mean by a term because of that difference.
\end{definition}

The theory of reference we rely on is that of LNG-02: a term refers to what the causal component of its use determines, namely the object or kind at the origin of an intact chain of reference-preserving uses grounded under a sortal that singles it out, and, when the causal component fails, to what the descriptive component determines, namely what satisfies the body of information associated with the term. For a natural kind term the grounding is in samples, and the term applies to whatever bears the relevant sameness relation to them, sameness of underlying nature for a substance term. We distinguish three elements of a speaker's internal state that bear on the reference of a kind term.

\begin{definition}\label{d:profile}
The \emph{reference-fixing profile} of a kind term in a speaker's internal state consists of:
\begin{enumerate}[label=(\alph*)]
\item the \emph{sortal} under which the grounding samples are taken (a kind of substance, a kind of organism, a kind of artefact);
\item the speaker's \emph{deferential and sameness intentions}: to apply the term to what bears the sameness relation appropriate to the sortal to the samples, and to use it as those from whom she acquired it do;
\item the \emph{associated information}: the descriptions she would offer of the term's extension.
\end{enumerate}
\end{definition}

\begin{principle}[Twin Earth]\label{p:twin}
Oscar and his twin in 1750 are internal duplicates, and Oscar's \enquote{water} applies to H$_2$O and not to XYZ, while his twin's applies to XYZ and not to H$_2$O.
\end{principle}

\begin{principle}[Sortal variation]\label{p:sortal}
Two speakers grounding a term in the same samples, in the same environment, under different sortals, mean different things by it.
\end{principle}

Principle~\ref{p:twin} is Putnam's verdict. Principle~\ref{p:sortal} is the lesson of the qua problem \citep{DevittSterelny1987}. A speaker who introduces \enquote{wug} pointing at a glass of water, intending it to apply to whatever is the same substance as this, means water by it. A speaker who introduces \enquote{wug} pointing at the same glass, intending it to apply to whatever is the same beverage as this, which is a functional sortal, means something whose extension includes XYZ in counterfactual situations where XYZ is drunk in the same way, and excludes the heavy water that the first speaker's term includes. A speaker whose intention is to apply the term to whatever has the same colour as this means something else again. The samples, the environment and the community are the same; the internal states differ; and so do the extensions.

\section{The argument}

\begin{theorem}\label{t:nothead}
Meanings are not in the head.
\end{theorem}

\begin{proof}
By Principle~\ref{p:twin}, Oscar and his twin are internal duplicates whose terms differ in extension. By Definition~\ref{d:meaning}, a difference in extension is a difference in meaning. So internal duplicates can differ in meaning, and by Definition~\ref{d:inhead} meanings are not in the head.
\end{proof}

\begin{theorem}\label{t:depend}
Meanings depend on the head.
\end{theorem}

\begin{proof}
By Principle~\ref{p:sortal}, two speakers in the same environment and community, grounding a term in the same samples under different sortals, differ in what they mean by it. The sortal is part of the speaker's internal state, so the speakers differ in internal state, and the difference in meaning is due to it, since nothing else differs. By Definition~\ref{d:inhead}, meanings depend on the head.
\end{proof}

\begin{proposition}\label{p:joint}
On the theory of reference of LNG-02, the meaning of a kind term is fixed by the speaker's internal state together with her environment and community: speakers who agree in all three mean the same by the term.
\end{proposition}

\begin{proof}
The causal component of a use is fixed by the grounding samples, which belong to the environment; by the sortal and the sameness intention, which belong to the internal state; and by the chain of acquisition, which belongs to the community and is engaged by the speaker's deferential intention. Whether the chain is intact is fixed by the same facts together with the causal history of the speaker's information. The descriptive component is fixed by the associated information, which is internal, and by what satisfies it, which is environmental. The hybrid rule of LNG-02 fixes reference as a function of the two components. Speakers who agree in internal state, environment and community therefore agree in both components and in reference, in the actual world; and since the sameness relation fixes the extension in counterfactual situations, they agree in meaning.
\end{proof}

Theorem~\ref{t:nothead} and Theorem~\ref{t:depend} are the two halves of the thesis. Proposition~\ref{p:joint} says that the two sources exhaust what fixes meaning: once the head and the world are both given, nothing further is needed. Burge's arthritis case shows that the world includes the community; the deferential intention is the point at which the community enters, and it is itself in the head.

\section{Dependence without description}

The thesis that meanings depend on the head must be reconciled with a thesis about kind terms on which MET-27-a relies: that kind terms refer independently of the descriptions associated with them. If the head contributes to meaning, and descriptions are in the head, it may seem that descriptions contribute to meaning after all.

The reconciliation lies in Definition~\ref{d:profile}. Of the three elements of the reference-fixing profile, only the third consists of descriptions of the extension, and it plays no role in fixing reference when the grounding succeeds. The first two are not descriptions that the referent must satisfy. The sortal says what kind of sameness relation to the samples is relevant, not what the samples are like: a speaker who grounds \enquote{water} under the sortal of a substance does not thereby believe anything about water's appearance, taste or location. The sameness and deferential intentions say how the term is to be extended from the samples and from the speaker's sources. A speaker can have them while believing falsehoods about every member of the extension, as the speaker who believes that water is an element does. Kripke's arguments from ignorance and error apply to the third element and not to the first two \citep{Kripke1980}.

The third element is not idle in every case. By the theory of LNG-02, it fixes reference when the causal component fails, either because a chain has broken or because a grounding does not single out one kind. \enquote{Jade} is the standard case of the second kind: its grounding samples were of two minerals with different compositions, and the term applies to whatever satisfies the associated information, an ornamental green stone of a certain hardness. Here, and only here, the associated descriptions fix the extension of a putative kind term, and the result is that the term does not refer to a natural kind at all. This is the restriction under which the independence thesis holds: kind terms whose groundings succeed refer independently of associated descriptions, and the descriptive fallback is engaged only when the causal chain breaks or fails to determine a referent. The two theses are therefore consistent, and each explains the scope of the other.

The same structure explains why the dependence of meaning on the head is easy to miss. In the cases on which the externalist argument is built, the internal states that differ between speakers are the associated descriptions, and these are idle; so internal variation makes no difference to meaning, and the head appears to contribute nothing. In the cases that show dependence, what varies is the sortal or the sameness intention, and these are rarely varied in the literature because they are shared by almost all users of a common kind term. Putnam's own account of the meaning of \enquote{water} included, besides the stereotype and the extension, semantic markers such as \emph{natural kind} and \emph{liquid} \citep{Putnam1975}. These are close to what we have called the sortal; our thesis is that they are in the head and that they fix meaning together with the samples.

\section{Objections and replies}

\begin{objection}[Narrow content]
Oscar and his twin share something semantic: a narrow content, which is a function from contexts to extensions and which explains their shared behavior \citep{Fodor1987}. This is the real meaning, and it is in the head.
\end{objection}

\begin{reply}
We accept that there is such a function; it is the contribution of the head that Theorem~\ref{t:depend} establishes. It is a function from environments to extensions, like Kaplan's character for indexicals, which is a function from contexts to contents \citep{Kaplan1989}. But it is not a meaning in the sense of Definition~\ref{d:meaning}, since it does not determine an extension; it determines one only given an environment. Calling it meaning abandons the role of meaning in fixing truth conditions, which is the role that made the question interesting. \citet{Segal2000} argues that narrow content is all psychology needs; we need not dispute this, since the thesis concerns meaning, not the taxonomy of psychological explanation.
\end{reply}

\begin{objection}[Two-dimensionalism]
The primary intension of \enquote{water} is the condition \enquote{the watery stuff of our acquaintance}, which is in the head and which determines, in each world considered as actual, what the term refers to \citep{Chalmers2002,Jackson1998}. Its extension-determining meaning is therefore in the head after all, in the form of a primary intension.
\end{objection}

\begin{reply}
If the primary intension is a description of watery appearance, it fails the arguments from error: water that is frozen, or heavily polluted, or that someone has never seen, is still water, and a speaker who has false beliefs about what water looks like still refers to it. If the primary intension is the condition of being the same substance as the samples that grounded the speaker's use, it is the causal component of LNG-02 written as a description, and it determines nothing until the samples are given. Either way, the intension is fixed in the head and the extension is fixed by the head and the world together. This is our thesis, not an objection to it.
\end{reply}

\begin{objection}[Self-knowledge]
A speaker knows a priori that she is thinking that water is wet. If externalism is true, she could deduce a priori that her environment contains water, which she cannot know a priori \citep{McKinsey1991}.
\end{objection}

\begin{reply}
The deduction requires the premise that having the concept water entails that there is water. On our account this premise is false. A speaker on a planet with no water, whose term \enquote{water} was introduced by a grounding that failed, has a term whose reference is fixed by its associated information, if anything satisfies it, or is empty \citep{Boghossian1997}. In neither case does she lack a concept; she has a different one, or an empty one. The speaker's a priori knowledge of what she is thinking is therefore compatible with her not knowing a priori what her environment contains, and the content of her self-ascription is fixed by the same facts that fix the content of the first-order thought, as \citet{Burge1988} argued.
\end{reply}

\begin{objection}[Social externalism]
Burge's case shows that what a speaker means depends on her community. The head drops out: the patient means arthritis whatever he has in mind.
\end{objection}

\begin{reply}
The patient means arthritis because he defers to his community's use, and the deference is part of his internal state. A speaker who used \enquote{arthritis} with the explicit intention of applying it to whatever ailment he has in his thigh, without deference, would mean something else, in the same community. Burge's case therefore shows that the community is part of what the head's contribution is applied to, not that the head contributes nothing.
\end{reply}

\begin{objection}[Extended minds]
If cognitive processes extend into reliably coupled external resources, as MND-09 argues, the boundary of the head is not the boundary of the internal state, and the question whether meanings are in the head is ill-formed.
\end{objection}

\begin{reply}
The externalism of MND-09 concerns the vehicles of cognitive processes; the externalism of Twin Earth concerns the determinants of content \citep{ClarkChalmers1998}. They are independent. If Oscar's notebook is part of his cognitive system, his twin has a duplicate notebook, and the case goes through with the extended system in place of the head. Definition~\ref{d:internal} may be read with the cognitive system in place of the body without affecting either theorem.
\end{reply}

\section{What the result does not establish}

The argument does not show that every term has a reference-fixing profile of the kind described in Definition~\ref{d:profile}; it shows this for kind terms and, through LNG-02, for names. It does not show that there is no useful notion of meaning on which meanings are in the head; it shows that the extension-determining notion is not one. And it does not decide how fine-grained the sortal must be for a grounding to succeed, which is a question about the qua problem that we have left with the theory of reference. The companion LNG-08-b examines variations of the Twin Earth case.

\section{Conclusion}

Meanings are not in the head: internal duplicates can mean different things, because the extensions of their terms depend on the samples in which the terms were grounded. Meanings depend on the head: speakers who share their samples, environment and community can mean different things, because the sortals and intentions with which they ground and extend their terms differ. What the head contributes is not a description of the extension, which is why kind terms refer independently of associated descriptions; the descriptions fix reference only when the grounding fails.

\begin{thebibliography}{99}\small
\bibitem[Boghossian(1997)]{Boghossian1997} Boghossian, P.~A. (1997). What the externalist can know a priori. \emph{Proceedings of the Aristotelian Society}, 97(2), 161--175.
\bibitem[Burge(1979)]{Burge1979} Burge, T. (1979). Individualism and the mental. \emph{Midwest Studies in Philosophy}, 4, 73--121.
\bibitem[Burge(1988)]{Burge1988} Burge, T. (1988). Individualism and self-knowledge. \emph{The Journal of Philosophy}, 85(11), 649--663.
\bibitem[Chalmers(2002)]{Chalmers2002} Chalmers, D.~J. (2002). The components of content. In D.~J. Chalmers (ed.), \emph{Philosophy of Mind: Classical and Contemporary Readings}, 608--633. Oxford University Press.
\bibitem[Clark and Chalmers(1998)]{ClarkChalmers1998} Clark, A. and Chalmers, D.~J. (1998). The extended mind. \emph{Analysis}, 58(1), 7--19.
\bibitem[Devitt and Sterelny(1987)]{DevittSterelny1987} Devitt, M. and Sterelny, K. (1987). \emph{Language and Reality: An Introduction to the Philosophy of Language}. MIT Press.
\bibitem[Fodor(1987)]{Fodor1987} Fodor, J.~A. (1987). \emph{Psychosemantics: The Problem of Meaning in the Philosophy of Mind}. MIT Press.
\bibitem[Jackson(1998)]{Jackson1998} Jackson, F. (1998). \emph{From Metaphysics to Ethics: A Defence of Conceptual Analysis}. Oxford University Press.
\bibitem[Kaplan(1989)]{Kaplan1989} Kaplan, D. (1989). Demonstratives. In J. Almog, J. Perry and H. Wettstein (eds.), \emph{Themes from Kaplan}, 481--563. Oxford University Press.
\bibitem[Kripke(1980)]{Kripke1980} Kripke, S.~A. (1980). \emph{Naming and Necessity}. Harvard University Press.
\bibitem[McKinsey(1991)]{McKinsey1991} McKinsey, M. (1991). Anti-individualism and privileged access. \emph{Analysis}, 51(1), 9--16.
\bibitem[Putnam(1975)]{Putnam1975} Putnam, H. (1975). The meaning of \enquote{meaning}. \emph{Minnesota Studies in the Philosophy of Science}, 7, 131--193.
\bibitem[Segal(2000)]{Segal2000} Segal, G. M.~A. (2000). \emph{A Slim Book about Narrow Content}. MIT Press.
\end{thebibliography}

\end{document}
